\documentclass{article}

\usepackage{amsmath,amssymb}
\usepackage{xurl}
\usepackage[hidelinks]{hyperref}
\setlength{\emergencystretch}{2em}

% Shared commands for notes in docs/paper-gaps/.

\DeclareUnicodeCharacter{2080}{\ifmmode{}_{0}\else\textsubscript{0}\fi}
\DeclareUnicodeCharacter{2081}{\ifmmode{}_{1}\else\textsubscript{1}\fi}
\DeclareUnicodeCharacter{2082}{\ifmmode{}_{2}\else\textsubscript{2}\fi}
\DeclareUnicodeCharacter{2099}{\ifmmode{}_{n}\else\textsubscript{n}\fi}

\newcommand{\Lean}{\textsc{Lean}}
\ExplSyntaxOn
\NewDocumentCommand{\leanid}{m}{
  \ifmmode
    \textsf{\detokenize{#1}}
  \else
    \group_begin:
    \urlstyle{sf}
    \tl_set:Nn \l_tmpa_tl {#1}
    \tl_replace_all:Nnn \l_tmpa_tl {\_} {_}
    \exp_args:NV \path \l_tmpa_tl
    \group_end:
  \fi
}
\ExplSyntaxOff
\providecommand{\tr}{\operatorname{tr}}
\newcommand{\tnleanIssueBase}{https://github.com/LionSR/TNLean/issues/}
\newcommand{\tnleanPrBase}{https://github.com/LionSR/TNLean/pull/}
\newcommand{\tnleanDocBase}{https://sirui-lu.com/TNLean/docs/}
\newcommand{\ghissue}[1]{\href{\tnleanIssueBase#1}{GitHub issue \##1}}
\newcommand{\ghpr}[1]{\href{\tnleanPrBase#1}{PR \##1}}
\newcommand{\doclink}[1]{\href{\tnleanDocBase#1.html}{\path{#1}}}

% Dirac notation and the identity operator, as in blueprint/src/macros/common.tex.
% \providecommand keeps note-local definitions authoritative.
\usepackage{dsfont}
\providecommand{\ket}[1]{|#1\rangle}
\providecommand{\bra}[1]{\langle #1|}
\providecommand{\braket}[2]{\langle #1 | #2 \rangle}
\providecommand{\ketbra}[2]{|#1\rangle\!\langle #2|}
\providecommand{\Id}{\mathds{1}}
\providecommand{\one}{\mathds{1}}
\providecommand{\C}{\mathbb{C}}
\providecommand{\MN}[1]{M_{#1}(\C)}
\providecommand{\MD}{M_D(\C)}

% External referencing for paper sources.  The build helper writes sanitized
% auxiliary files under build/paper-aux/ and excludes duplicated source labels.
\usepackage{xr}

% Import a generated paper auxiliary file with a caller-supplied label prefix.
% The candidate paths support builds from docs/paper-gaps/, the repository root,
% and build/paper-aux/ itself.
\newcommand{\importpaperaux}[2]{%
  \IfFileExists{../../build/paper-aux/#2.aux}{%
    \externaldocument[#1]{../../build/paper-aux/#2}%
  }{%
    \IfFileExists{build/paper-aux/#2.aux}{%
      \externaldocument[#1]{build/paper-aux/#2}%
    }{%
      \IfFileExists{#2.aux}{%
        \externaldocument[#1]{#2}%
      }{}%
    }%
  }%
}

% General external reference macros
\newcommand{\paperref}[2]{\ref{#1-#2}}
\newcommand{\papereqref}[2]{(\ref{#1-#2})}

% CPSV16 source (1606.00608 / MPDO-22-12-17-2.tex) external document and shortcuts
\importpaperaux{cpsv16-}{1606.00608/MPDO-22-12-17-2-tnlean}
\newcommand{\cpsvref}[1]{\paperref{cpsv16}{#1}}
\newcommand{\cpsveqref}[1]{\papereqref{cpsv16}{#1}}

% Verdict marker: \gapnote{<kind>}{<status>}. Kind is the policy
% classification (clarification, local-correction, scope-restriction,
% unfaithful, false-source, open-gap); status is open, wip (elimination
% underway), resolved, or historical. Machine-read by the site index;
% prints a verdict line.
\newcommand{\gapnote}[2]{\par\noindent\textbf{Verdict:} #1 (#2).\par}


\title{The One-Step Subspace in the Quantum Wielandt Inequality}
\author{}
\date{2026-04-30}

\begin{document}
\maketitle

\section{The assertion}

The quantum Wielandt inequality of Sanz--P\'erez-Garc\'ia--Wolf--Cirac is
stated in \cite[Theorem~1]{Sanz2010Quantum}.\footnote{For
    traceability, this note corresponds to
    \href{https://github.com/LionSR/TNLean/issues/1049}{issue \#1049}. The local
    source passage is \path{Papers/0909.5347/main.tex:360--365},
    \path{Papers/0909.5347/main.tex:572--590}, and
    \path{Papers/0909.5347/main.tex:645--684}.}
Let $\mathcal E_A$ be a quantum channel on $M_D(\mathbb C)$ with Kraus
operators $A_1,\ldots,A_a$. For each $n$, set
\[
    S_n(A)=\operatorname{span}\{A_{i_1}\cdots A_{i_n}:1\leq i_1,\ldots,i_n\leq a\}.
\]
Thus $S_n(A)$ is the subspace generated by words of length exactly $n$. Let
\[
    i(A)=\min\{n:S_n(A)=M_D(\mathbb C)\},
\]
when this minimum exists.

The paper writes $d$ for the number of linearly independent Kraus operators.
Equivalently, after replacing the displayed Kraus list by a minimal spanning
list, one has
\[
    d=\dim S_1(A).
\]
With this convention, Theorem~1 asserts that a primitive channel satisfies
\[
    i(A)\leq (D^2-d+1)D^2.
\]
The same theorem also gives two better estimates under additional hypotheses
on the one-step subspace. If $S_1(A)$ contains an invertible matrix, then
\[
    i(A)\leq D^2-d+1.
\]
If $S_1(A)$ contains a non-invertible matrix with a nonzero eigenvalue, then
\[
    i(A)\leq D^2.
\]
These are hypotheses about arbitrary matrices in $S_1(A)$. The paper does not
require the special matrix to be one of the Kraus operators in a chosen Kraus
representation.

\section{What has been formalized}

The formal statements use the same subspaces $S_n(A)$ and define the Kraus
rank by
\[
    \operatorname{kr}(A)=\dim S_1(A).
\]
This is the paper's number $d$ written invariantly. Replacing $d$ by
$\operatorname{kr}(A)$ is therefore not a change in the mathematical bound.
\footnote{The formal definitions are \leanid{MPSTensor.wordSpan} in
    \path{TNLean/Wielandt/SpanGrowth/CumulativeSpan.lean} and
    \leanid{MPSTensor.krausRank} in
    \path{TNLean/Wielandt/Primitivity/Definitions.lean}. The local
    source identifies $d$ as the number of linearly independent Kraus operators in
    \path{Papers/0909.5347/main.tex:360--365}; the proof of Lemma~1 uses
    $\dim T_1(A)=d$ in \path{Papers/0909.5347/main.tex:579--583}.}

The general inequality is formalized for the same index $i(A)$:
\[
    i(A)\leq (D^2-\operatorname{kr}(A)+1)D^2.
\]
The proof passes through the usual eventual full-rank formulation and a
blocking argument, but the conclusion is the equality
$S_n(A)=M_D(\mathbb C)$ for one fixed length $n$. It is not merely a statement
about the increasing sum $S_0(A)+\cdots+S_n(A)$.\footnote{The general
    theorem is \leanid{MPSTensor.iIndex_le_general_of_isPrimitivePaper}
    in \path{TNLean/Wielandt/Inequality/Bounds.lean}. The formal equivalence
    between the primitive condition used by Sanz--P\'erez-Garc\'ia--Wolf--Cirac
    and eventual full Kraus rank under normalization is in
    \path{TNLean/Wielandt/Primitivity/Equivalence.lean}.}

The two improved estimates now use the source hypotheses directly. If an
arbitrary $X\in S_1(A)$ is invertible, then
\[
    S_{D^2-\operatorname{kr}(A)+1}(A)=M_D(\mathbb C).
\]
If an arbitrary $X\in S_1(A)$ is non-invertible and has a nonzero eigenvalue,
then
\[
    S_{D^2}(A)=M_D(\mathbb C).
\]
Thus the distinguished-Kraus-operator restriction recorded in the first
version of this note has been removed.\footnote{The primary declarations are
    \leanid{MPSTensor.wordSpan_eq_top_of_isPrimitivePaper_of_mem_wordSpan_one_of_isUnit}
    and
    \leanid{MPSTensor.wordSpan_eq_top_of_isPrimitivePaper_of_mem_wordSpan_one_of_noninvertible_eigenvector}
    in \path{TNLean/Wielandt/Inequality/Bounds.lean}. The corresponding index
    bounds are \leanid{MPSTensor.iIndex_le_of_mem_wordSpan_one_of_isUnit} and
    \leanid{MPSTensor.iIndex_le_sq_of_mem_wordSpan_one_of_noninvertible_eigenvector}.}

\section{Resolution of the one-step-subspace deviation}

The original restriction arose from proving the dimension-growth and rank-one
arguments only for a distinguished Kraus matrix. The current statements take an
arbitrary $X\in S_1(A)$ together with its span-membership hypothesis. The key
inclusion
\[
    XS_n(A)\subseteq S_{n+1}(A)
\]
then replaces the distinguished-letter step. For the invertible case, left
multiplication by $X$ is injective and drives dimension growth. For the
non-invertible case, the same inclusion pads the rank-one operator extracted
from a nonzero eigenvector. Thus both proofs now apply to exactly the one-step
subspace hypotheses used in the source, and the deviation recorded by the first
version of this note is closed.

\section{What is not a discrepancy}

The replacement of $d$ by $\operatorname{kr}(A)=\dim S_1(A)$ is not a paper
error and not a formal weakening. It records explicitly the linearly independent
Kraus count used by the source proof.

Nor is the use of the increasing spaces
\[
    T_n(A)=S_0(A)+S_1(A)+\cdots+S_n(A)
\]
by itself a theorem-level deviation. Some intermediate formal results prove
$T_{D^2}(A)=M_D(\mathbb C)$ under normality. These results are weaker than the
full quantum Wielandt theorem, but the theorem in
\path{TNLean/Wielandt/Inequality/Bounds.lean} proves the stated bound for
$i(A)$, where $i(A)$ is defined using the spaces $S_n(A)$ themselves.
\footnote{The cumulative-span theorem is
    \leanid{MPSTensor.cumulativeSpan_eq_top} in
    \path{TNLean/Wielandt/SpanGrowth/NonzeroTraceProduct.lean}; the theorem for
    the spaces $S_n(A)$ is the general theorem cited above. The blueprint makes
    this distinction at \texttt{thm:cumulative_eq_top} and states the general
    bound at \texttt{thm:wielandt_general} in the Chapter~8 sources under
    \path{blueprint/src/chapter/}.}

Finally, the lower-level use of normality is not itself a mismatch. In these
formal statements, normality is eventual full Kraus rank. Under the paper's
normalization, the formal statements prove the equivalence between the paper's
primitive condition and eventual full Kraus rank, so using normality as an
intermediate hypothesis is legitimate.

\section{Scope of the index comparison}

For a nonzero vector $\varphi\in\mathbb C^D$, set
\[
    H_n(A,\varphi)=S_n(A)\varphi,
\]
and let $q(\mathcal E_A)$ be the least positive $n$ such that
$H_n(A,\varphi)=\mathbb C^D$ for every nonzero $\varphi$. As for $i(A)$, the
value is defined to be zero when no such positive integer exists.
Proposition~1 of~\cite{Sanz2010Quantum} states
\[
    q(\mathcal E_A)\leq i(A)
\]
for every quantum channel.

The available comparison assumes trace-preserving normalization and uniform
vector-spreading primitivity. It proves the same inequality under these
hypotheses, and is therefore a specialization of Proposition~1 rather than the
source theorem with its full hypothesis set. A related comparison assumes
instead that $i(A)$ is finite.\footnote{The two declarations are
    \leanid{MPSTensor.qIndex_le_iIndex_of_isPrimitivePaper} and
    \leanid{MPSTensor.qIndex_le_iIndex}. The definitions are in
    \path{TNLean/Wielandt/Primitivity/Definitions.lean}, and the restricted
    theorem is in \path{TNLean/Wielandt/Inequality/Bounds.lean}.}

The missing cases are mathematically degenerate but must still be represented
explicitly. If the channel is not primitive, then the admissible set defining
$q(\mathcal E_A)$ is empty and $q(\mathcal E_A)=0$. If it does not have
eventual full Kraus rank, the admissible set for $i(A)$ is empty and $i(A)=0$.
A source-faithful theorem can therefore be obtained by proving these two
zero-index characterizations and splitting on primitivity, using the existing
primitive-channel inequality in the nontrivial case.

\section{Scope of the strong-irreducibility equivalence}

Proposition~3(c) of~\cite{Sanz2010Quantum} asks for a unique peripheral
eigenvalue and a positive-definite corresponding eigenvector. The present
comparison strengthens this condition by also requiring irreducibility of the
transfer map. Consequently, the proved equivalence is a strengthened version
of Proposition~3 rather than its literal third clause.\footnote{The strengthened
    predicate is \leanid{MPSTensor.IsStronglyIrreduciblePaper}; the equivalences
    are \leanid{MPSTensor.primitivePaper_iff_stronglyIrreducible} and
    \leanid{MPSTensor.hasEventuallyFullKrausRank_iff_stronglyIrreducible}. See
    \path{TNLean/Wielandt/Primitivity/Definitions.lean} and
    \path{docs/glossary.md}.}

To remove the restriction, define the literal source predicate using only the
positive-definite fixed point and peripheral-eigenvalue uniqueness, then prove
that it implies transfer-map irreducibility for a normalized completely
positive map. This implication supplies the extra conjunct of the existing
predicate, after which the established equivalences yield the literal
Proposition~3. Alternatively, the three equivalence directions may be reproved
directly for the source predicate.

\section{Conclusion}

The general quantum Wielandt bound and both improved one-step-subspace estimates
now use the source indices and hypotheses. Two local scope restrictions remain:
the index comparison assumes primitivity, and the strong-irreducibility
predicate carries an additional irreducibility conjunct. The preceding sections
record the corresponding zero-index and source-predicate bridges needed to
remove these restrictions.

\bibliographystyle{alpha}
\bibliography{references}

\end{document}
